\documentclass[11pt,letterpaper]{article}
\usepackage[margin=1in]{geometry}\usepackage{amsmath,amssymb,booktabs,array,microtype}\usepackage[T1]{fontenc}\usepackage{lmodern}
\usepackage[colorlinks=true,linkcolor=blue,citecolor=blue,urlcolor=blue]{hyperref}
\setlength{\parindent}{0pt}\setlength{\parskip}{5pt}\sloppy
\title{The Decoupling Locus:\\ a MOND Scale Derived from a Clock,\\ and the Solar System Reduced to One Number}
\author{Carl P.\ Zimmerman\\ \small Briar Creek Tech\\ \small AI-assisted research programme; not peer reviewed}
\date{12 September 2026}
\begin{document}\maketitle
\begin{abstract}
A cuscuton-clock sector with a projected-gradient scalar produces a cold clustering component and passes nine cosmological and solar-system gates as a dark component \cite{paper19,paper20,paper21,paper22}, but produces no force law, because matter in that action is minimally coupled: $\delta S_{\rm matter}/\delta\chi=0$ identically. Adding a coupling is what would give MOND, and the obstruction was that the clock and the scalar mix, so a coupled matter source would drag the clock and generate preferred-frame effects. We report that the mixing vanishes \emph{exactly} on the locus $W_0\equiv U-2\gamma\bar q^2\dot{\bar q}=0$, which is a condition on coefficients the action already has. A first pass priced that locus at a cubic coupling eight orders too large, scaling as $1/w$; that pricing held the clock coefficient $U$ fixed, and $U$ is not free. With $\rho=U/m_{\rm rel}$ imposed the factor of $w$ cancels identically: what the locus fixes is the dimensionless ratio $\gamma\bar q^3H/\rho=-1/[6(s_0-1)]$, of magnitude at most $1/6$, and the operator's weight against gravity is $\Omega_{\rm sec}/[6(s_0-1)]\le0.044$. The earlier verdict is withdrawn. The locus also \emph{forces} two previously free quantities: positivity of the sector's remaining energy requires $s_0\ge2$, hence $m_{\rm rel}\le w\lesssim10^{-4}$. Extending the action by the $Y^{3/2}$ operator identified as missing in \cite{paper23} and a leading-order matter coupling, the quasi-static gradient sector yields the deep-MOND scaling outright ($d\log g/d\log r=-1$, $d\log g/d\log M=+1/2$) and fixes the acceleration scale, $a_0=\lambda^3/(12\pi G\beta s_0)$, mass-independent and falling with the clock rate. The single conflict --- the branch wants the linear gradient coefficient large, deep MOND wants it small, and it is the same coefficient --- resolves into one inequality, $C\equiv2k\bar q^2/U\gtrsim10^3$, whose derived exponents cancel exactly, so $C$ is a constant of the motion. Finally, coupling matter \emph{disformally} along the clock's own unit timelike direction shifts both weak-field potentials equally, giving $\gamma_{\rm PPN}=1$ exactly with no added vector field; boosted, the same term gives $\alpha_1=8f_s$, five orders over the bound unless local matter drags the clock to within $4.6\,$m$\,$s$^{-1}$. That alignment, and $\kappa$, are what remain open. Ten new Lean~4 certificates; the file stands at 124 theorems, zero \texttt{sorry}.
\end{abstract}

\section{Position, and three corrections to the previous paper}
The action \cite{paper13} is, with $c=1$ and signature $(-{+}{+}{+})$,
\begin{equation}
S=\int d^4x\sqrt{-g}\Big[\tfrac{M^2}{2}(R-2\Lambda)+P(X,\tau)-V(\tau)+sW(Y,\tau)+\gamma X\Box\chi\Big]+S_r+S_b,
\label{action}
\end{equation}
with $s=\sqrt{-\nabla_\mu\tau\nabla^\mu\tau}$, $n_\mu=-\partial_\mu\tau/s$, $Q=n^\mu\partial_\mu\chi$, $Y=(g^{\mu\nu}+n^\mu n^\nu)\partial_\mu\chi\partial_\nu\chi$, $X=Q^2-Y$. An external audit of \cite{paper23} found three errors, each reproduced independently before being accepted, and they are recorded here because two of them were load-bearing.

\emph{(i) The health condition had the wrong sign.} We stated $4d\ell>U$. Checking the same reduction against the canonical scalar $P(X)=X/2$, which gives $F_Y=-\tfrac12$, shows a healthy field has \emph{negative} $F_Y$. The correct condition is $\boxed{U>4d\ell}$. Every verdict resting on the reversed inequality is withdrawn, including a proposed coefficient window.

\emph{(ii) The interpolating exponent was measured against the wrong acceleration.} The reported slope $0.470$ compared the kernel with the Newtonian acceleration $x=g_N/a_0$ rather than the true one $y=x\nu(x)$. Converted, the slope is $0.997$. The qualitative conclusion of \cite{paper23} --- that MOND requires a linear $\mu$ in $y$, which the unextended action does not supply --- survives with the number corrected.

\emph{(iii) The no-force-law verdict was right for a different reason.} \cite{paper23} excluded a force law on the exponents of the two gradient operators. The stronger and simpler fact is that matter in (\ref{action}) enters only through $S_r$ and $S_b$, both minimally coupled to $g_{\mu\nu}$: $\delta S_{\rm matter}/\delta\chi=0$ identically. The action modifies \emph{no} force law, and the rotation-curve exponents of \cite{paper23} are conditional on a coupling that is not present. Correspondingly, a structural argument that a vanishing background gradient removes the clock--scalar mixing is also withdrawn: the exact quadratic action carries a cross term $-2qW_Y$ and a clock response $\sigma/\pi=2qsW_Y/(2q^2W_Y-W_0)$, both of which survive $Y\to0$. The preferred-frame gate returns to open, and that is the gate this paper addresses.

\section{The locus on which the clock stops listening}
The clock response above vanishes identically, together with the entire $W_Y$ contribution to the eliminated scalar coefficient, when
\begin{equation}
W_0\;\equiv\;U-2\gamma\bar q^2\dot{\bar q}\;=\;0 ,
\label{locus}
\end{equation}
at which point that coefficient reduces to $2P_X$ alone. Since $W_0$ is built from quantities (\ref{action}) already contains, this is a locus inside the existing parameter space, not a modification. It has one immediate cost: with $W_0=0$ every surviving term of the clock equation carries the same factor of the clock rate, which therefore cancels. That equation no longer determines $s_0$ but constrains the coefficient functions; $s_0$ remains fixed by conservation through $s_0-1=w/m_{\rm rel}$.

\section{What the locus costs, measured in units that cannot be inflated}
\emph{The stress tensor, derived rather than quoted.} Write $S=\int dt\,Na^3L_m$ in minisuperspace with the lapse restored. Then $\rho=-a^{-3}\partial_N(Na^3L_m)|_{N=1}$ and $p=(3a^2)^{-1}[\partial_a-d_t\partial_{\dot a}](Na^3L_m)|_{N=1}$. Calibrating on $L_m=F(Q)$ reproduces the textbook pair $QF_Q-F$ and $F$. Applied to the cubic operator, whose minisuperspace action is $-2\gamma a^2\dot a Q^3$ after an integration by parts that differs from the unintegrated form by an explicit total derivative,
\begin{equation}
\rho_3=-6\gamma H\bar q^3,\qquad p_3=2\gamma\bar q^2\dot{\bar q},
\end{equation}
so that $W_0=U-p_3$: \emph{the decoupling condition says the cubic operator's pressure equals the clock coefficient}. On the derived family $\bar q\propto a^{-3w}$ one finds $p_3=w\rho_3$ identically, so the operator carries the sector's own equation of state and switching it on does not disturb the background scaling.

\emph{The correction.} Using the family's $\rho=U/m_{\rm rel}$ and $p=U(s_0-1)$ \cite{paper21}, the cubic operator's share of the sector density is
\begin{equation}
\frac{\rho_3}{\rho}=\frac{1}{s_0-1},
\end{equation}
free of both $U$ and $w$. A first pass had reported $|\gamma|=U/(6|w|\bar q^3H)$ and concluded, holding $U$ fixed, that the acoustic bound $w\lesssim10^{-4}$ forces $|\gamma|\sim650$ against the $10^{-6}$ used throughout. But $U=m_{\rm rel}\rho$ with $m_{\rm rel}=w/(s_0-1)$, and substituting, $\partial|\gamma|/\partial w=0$ exactly. The divergence was an artifact of units in which the sector density had been set to $1/m_{\rm rel}$, which is large for the same reason $w$ is small. What the locus fixes is a ratio,
\begin{equation}
\frac{\gamma\bar q^3H}{\rho}=-\frac{1}{6(s_0-1)},
\end{equation}
of magnitude at most $1/6$; today $|\gamma|\bar q^3=1.13\times10^{21}\,$eV$^3$, a cube root of $10.4\,$MeV. Comparing the cubic operator's share of the momentum conjugate to the scale factor with the Einstein--Hilbert term's, from the same action,
\begin{equation}
\left|\frac{\pi_3}{\pi_{\rm grav}}\right|=\frac{\Omega_{\rm sec}}{6(s_0-1)}\;\le\;0.044 ,
\end{equation}
a percent-level correction. Every result derived at $\gamma\to0$ therefore survives on this locus to about four percent, and the earlier verdict that they were invalidated is withdrawn.

\emph{Two free quantities become forced.} The non-cubic density is $\rho[1-1/(s_0-1)]$, which is non-negative only for
\begin{equation}
s_0\ge2\qquad\Longleftrightarrow\qquad m_{\rm rel}\le w\lesssim10^{-4}.
\end{equation}
The clock must run at least twice proper time, and the margin must sit \emph{below} the equation of state rather than above it. At the boundary $s_0=2$ the cubic operator supplies the entire dark sector density on its own. Both statements are Lean-certified.

\section{The force law on the locus}
Extend $W$ by the $Y^{3/2}$ operator \cite{paper23} identified as the missing ingredient, and couple matter at leading order, $\lambda\varphi\rho_m$. In the quasi-static limit the gradient equation is $\nabla\!\cdot\!(2sW_Y\nabla\varphi)=-\lambda\rho$, so $2sW_Y$ plays the role of the interpolating function. With $W=\beta Y^{3/2}$ the spherical solution gives an acceleration on matter with
\begin{equation}
\frac{d\log g}{d\log r}=-1,\qquad \frac{d\log g}{d\log M}=+\tfrac12 ,
\end{equation}
the deep-MOND scaling, obtained from the action rather than assumed. Matching to $\sqrt{GMa_0}/r$ fixes the scale in terms of the action's own coefficients,
\begin{equation}
\boxed{\;a_0=\frac{\lambda^3}{12\pi G\beta s_0}\;}\qquad\text{with}\qquad \frac{\partial a_0}{\partial M}=0 ,
\end{equation}
so it is a constant of the theory rather than an object-by-object fit, and it \emph{falls as the clock runs faster}. On the canonical footing $\lambda^3/(\beta s_0)=2.3555\times10^{-19}$ in SI units and on the alternative footing $2.8378\times10^{-19}$, the ratio being the footing fork $1.204777$ itself.

\emph{The one conflict, and its window.} The linear term $d$ in $W_Y$ appears twice: as the floor $\mu_{\rm floor}=d/k$ it places under the interpolating function, $k$ being $W_Y$ at the highest acceleration probed, and as the closure's $\mu=2d\bar q^2/U$ that sets the clock margin. They are the same $d$, related by $\mu=C\,\mu_{\rm floor}$ with $C=2k\bar q^2/U$. The locus needs $\mu\ge1-w$; the rotation-curve data need $\mu_{\rm floor}\le10^{-3}$ for the radial acceleration relation to survive at the ten-percent level down to $g/a_0=10^{-2}$. Both hold exactly when
\begin{equation}
C\;=\;\frac{2k\bar q^2}{U}\;\gtrsim\;10^{3}.
\end{equation}
On the derived family the exponents cancel, $-3(1-w)-6w+3(1+w)=0$, so \emph{$C$ is a constant of the motion}: imposed once, it holds for all time. Missing the window does not kill the construction but buys it with a cancellation among the sector's densities of degree $m_{\rm rel}/w$, which reaches $10^4$ at $C=1$ --- the same $1/w$ the first pass saw, relocated from the coupling to the cancellation, where it belongs.

\section{The solar system, reduced to one number}
A \emph{conformal} matter coupling shifts the time potential and not the space one and fails light bending outright. The standard cure is a disformal coupling along a unit timelike vector, which TeVeS \cite{teves} and AeST \cite{aest} must add as a separate dynamical field. Action (\ref{action}) already has one: the clock's gradient supplies $n_\mu$, and the cuscuton \cite{cuscuton} propagates no scalar of its own. Expanding $\tilde g=(1-2\varphi)g-4\varphi\,n\otimes n$ in the weak field,
\begin{equation}
\tilde\Phi=\Phi+\varphi,\qquad \tilde\Psi=\Psi+\varphi\qquad\Longrightarrow\qquad \gamma_{\rm PPN}=1\ \text{exactly},
\end{equation}
however strong the scalar: the shift is common to both potentials and cancels from their difference. Light bending and Shapiro delay are general-relativistic.

Boosting is where it costs. For a system moving at $w$ through the clock's frame the same term gives $\tilde g_{0i}=4\varphi w_i$, first order in the velocity \emph{and} first order in the scalar, which is the preferred-frame signature. Matching to the post-Newtonian form,
\begin{equation}
\alpha_1=8f_s ,
\end{equation}
$f_s$ being the scalar's share of the local potential. At the share MOND requires at solar-system accelerations, $f_s\approx1$, that is $8$ against a bound of $10^{-4}$: over by $8\times10^4$. A clock sitting in the cosmic frame is excluded outright, by the same wall that closed the vector-based completions. The bound converts directly into an alignment requirement:
\begin{equation}
\boxed{\;v_{\rm residual}\le4.6\ \text{m\,s}^{-1}\;}
\end{equation}
between the clock's frame and the local matter frame, one part in $8\times10^4$ of the solar system's $370\,$km$\,$s$^{-1}$ through the cosmic frame. \emph{Whether the clock's own field equation delivers that is not computed here}, and an earlier lane measuring a cuscuton clock as dragged at $3\times10^{-2}\Phi_N$ at Saturn \cite{paper13} is the same question in a different parameterisation and was recorded there as a kill.

\section{The board, and what is not claimed}
Every gate the programme has posed now carries a computed number on this locus --- the preferred-frame mixing (identically zero), the cubic backreaction ($\le0.044$), the clock rate ($\ge2$) and margin ($\le w$), the acoustic scale, the forest through criticality \cite{paper20}, lensing and $S_8$ (the sector is inert), and the force law and its scale --- with two exceptions, which are stated as exceptions. The clock alignment above is not computed. And $\kappa=\tfrac12$ remains fitted and provably underivable by this class of actions \cite{repo}: that is a theorem of this programme, not a gap in this paper. \textbf{This is not a complete theory of gravity and nothing here should be read as one.}

Limits, in full. The elimination giving the mixing coefficient is an external audit's, used rather than re-derived. The positivity argument forcing $s_0\ge2$ assumes no piece of the sector carries negative energy; if one may, $s_0<2$ costs a cancellation of degree $(2-s_0)/(s_0-1)$. The weight $\pi_3/\pi_{\rm grav}$ is a background measure; the perturbation kinetic matrix at $s_0=2$ has not been recomputed and no gate has been re-run numerically on the locus. The matter coupling is leading order, not a full conformal or disformal factor; the quasi-static reduction drops the clock's gradient response, legitimate only because the mixing vanishes here; $k$ is defined, not derived; the interpolation used is the pure $Y^{3/2}$ limit, not $\nu_{\rm RAR}$ in full, so the deep-MOND limit and the scale are established and the whole interpolating function is not. The post-Newtonian matching uses the standard normalisation, so the coefficient $8$ is convention-dependent at the factor-of-two level; the five-orders conclusion is not.

\section{Reproducibility}
\texttt{fable\_independent\_2026/} lanes \texttt{L211\_three\_corrections.py}, \texttt{L212\_decoupling\_branch.py} (commit \texttt{473860435}), \texttt{L213\_branch\_solved\_and\_the\_board.py}, \texttt{L214\_force\_law\_on\_the\_branch.py} (commit \texttt{425a28355}) and \texttt{L215\_solar\_system\_on\_the\_branch.py} (commit \texttt{bde5bfd2c}), each with its committed output and results file. The Lean file \texttt{fable\_independent\_2026/lean\_2026/Mondlean.lean} carries \texttt{cubic\_share\_is\_clock\_rate}, \texttt{positivity\_forces\_clock\_rate}, \texttt{margin\_below\_equation\_of\_state}, \texttt{cubic\_coupling\_is\_equation\_of\_state\_independent}, \texttt{mond\_scale\_mass\_independent}, \texttt{branch\_mond\_window}, \texttt{derived\_exponents\_cancel}, \texttt{window\_ratio\_is\_scale\_invariant}, \texttt{disformal\_gamma\_is\_one} and \texttt{preferred\_frame\_alignment\_requirement}; 124 theorems, exit 0, zero \texttt{sorry}, axioms $\subseteq\{$\texttt{propext}, \texttt{Classical.choice}, \texttt{Quot.sound}$\}$.

\begin{thebibliography}{9}
\bibitem{repo} The research repository, \url{https://github.com/carlzimmerman/zimmerman-formula} (2026).
\bibitem{paper13} C.~P.~Zimmerman, \emph{A Primordial-Clock Relativistic MOND Candidate: Status Report}, doi:10.5281/zenodo.22699828 (2026).
\bibitem{paper19} C.~P.~Zimmerman, \emph{Gradient-Driven Criticality}, doi:10.5281/zenodo.22727111 (2026).
\bibitem{paper20} C.~P.~Zimmerman, \emph{The Coldness Is Dynamical}, doi:10.5281/zenodo.22727860 (2026).
\bibitem{paper21} C.~P.~Zimmerman, \emph{The Clock Rate Is the Dark Matter Pressure}, doi:10.5281/zenodo.22728545 (2026).
\bibitem{paper22} C.~P.~Zimmerman, \emph{A Clock That Cannot Be Dragged}, doi:10.5281/zenodo.22728789 (2026).
\bibitem{paper23} C.~P.~Zimmerman, \emph{What a Clock Can and Cannot Do}, doi:10.5281/zenodo.22728946 (2026).
\bibitem{bekenstein} J.~Bekenstein, M.~Milgrom, Astrophys.\ J.\ \textbf{286}, 7 (1984).
\bibitem{cuscuton} N.~Afshordi, D.~J.~H.~Chung, G.~Geshnizjani, Phys.\ Rev.\ D \textbf{75}, 083513 (2007).
\bibitem{teves} J.~D.~Bekenstein, Phys.\ Rev.\ D \textbf{70}, 083509 (2004).
\bibitem{aest} C.~Skordis, T.~Z\l{}o\'snik, Phys.\ Rev.\ Lett.\ \textbf{127}, 161302 (2021).
\end{thebibliography}
\end{document}
